\scanpage{12}

\begin{center}- 2 -\end{center}
\margin{$\mathrm{n}^{\mathrm{o}}398$}

PROPOSITION 2.-- \uline{Si} $E'\longrightarrow E\longrightarrow E''$ \uline{est une suite exacte de} $A$-\uline{modules de longueur finie, la suite} $E''^*\longrightarrow E^*\longrightarrow E'^*$ \uline{est exacte}.

Soit d'abord $0\longrightarrow L\longrightarrow M\longrightarrow N\longrightarrow0$ une suite exacte de modules de longueur finie et montrons que $0\longrightarrow N^*\longrightarrow M^*\longrightarrow L^*$ est exacte. C'est en tout cas vrai pour $0\longrightarrow N^*\longrightarrow M^*\longrightarrow L^*$, d'après les propriétés générales de $\operatorname{Hom}$. Comptant les longueurs, on en déduit que $M^*\longrightarrow L^*$ est surjectif, d'où le résultat.

Le cas d'une suite exacte qqe $E'\longrightarrow E\longrightarrow E''$ se ramène au précédent en dévissant à la façon habituelle, (Chap. II, 3e éd., p. 21). Soit $E$ un $A$-module. Si $F\subset E$, on note $F^{\circ}$ son orthogonal dans $E^*$, i.e. etc $\ldots$

PROPOSITION 3.-- \uline{Soit} $E$ \uline{un} $A$-\uline{module de longueur finie}. \uline{La correspondance qui à} $F\subset E$ \uline{fait correspondre} $F^{\circ}$ \uline{est une bijection de l'ensemble des sous-modules de} $E$ \uline{sur celui des sous-modules de} $E^*$ ; \uline{cette bijection transforme somme en intersection et intersection en somme}.

On utilise la suite exacte $0\longrightarrow(E/F)^*\longrightarrow E^*\longrightarrow F^*\longrightarrow0$, qui montre que $F^*$ s'identifie à $E^*/F^{\circ}$ et $(E/F)^*$ à $F^{\circ}$. On en déduit que l'injection $F\longrightarrow F^{\circ\circ}$ s'identifie à l'injection $F\longrightarrow F^{**}$, qui est un isomorphisme. D'où, etc $\ldots$

\uline{Exemple}.-- $A=$ \uuline{$Z$}, $K/A=$ \uuline{$Q$}$/$\uuline{$Z$} ; on a une dualité pour les groupes abéliens finis.

\begin{center}-:-:-:-\end{center}

\mspage{11}
